% === Experimental Setup ===
\section{Experimental Setup}
\label{sec:setup}

\subsection{Research Questions}
\label{sec:setup-rqs}

Our evaluation is organized around four questions:

\begin{itemize}
  \item[\textbf{RQ1}] \textbf{(Retrieval).} How accurately can a live incident be
  matched to the most similar resolved incident in memory, and does fusing dense,
  lexical, and graph retrieval outperform each in isolation, as well as strong
  prior-art and LLM-based retrievers?
  \item[\textbf{RQ2}] \textbf{(End-to-end triage).} Does the capability-gated
  system retrieve effectively end to end, and does the \emph{same} system hold up
  as the available evidence degrades from full telemetry to text alone?
  \item[\textbf{RQ3}] \textbf{(Cost of adaptive execution).} What does the system's
  capability gating buy? How much triage cost does it save relative
  to running every step on every window, at equal accuracy, and which components
  carry the result?
  \item[\textbf{RQ4}] \textbf{(Synthetic vs.\ real).} How do triage and retrieval
  behave on synthetic versus real incidents, and what does the contrast reveal
  about where the difficulty lies?
\end{itemize}

\subsection{Systems and Baselines}
\label{sec:setup-baselines}

We compare our hybrid retriever and full system against three groups of baselines,
all evaluated under an identical memory, time-aware visibility, and gold contract.

\noindent\emph{Constituent retrievers.} To show fusion exceeds the sum of its parts, we report our fine-tuned dense
bi-encoder and our knowledge-graph retriever on their own, alongside the full
hybrid that fuses them with a learned sparse retriever.

\noindent\emph{Prior-art retrieval.} We add widely used, off-the-shelf retrievers
applied zero-shot: three strong dense embedding models (BGE, E5, and
MPNet)~\cite{xiao2024cpack,wang2022e5,song2020mpnet}, classical lexical search
(TF--IDF and BM25~\cite{robertson2009bm25}), and a cross-encoder
re-ranker~\cite{nogueira2019passage} applied over the dense candidates.

\noindent\emph{Retrieval-augmented LLM.} As a contemporary comparison,
an instruction-tuned language model (Qwen2.5-7B~\cite{qwen2024}) re-ranks the
dense top candidates per query; following common practice, we run it on a fixed
500-incident sample per dataset.

\noindent\emph{Triage classifiers.} For the triage decision (RQ4) we additionally
compare standard classifiers over the numeric telemetry features (gradient-boosted
trees, random forests, and logistic regression) wherever such features exist.

\subsection{Metrics}
\label{sec:setup-metrics}

We measure two capabilities separately. \emph{Retrieval} is scored with
Hit@$k$ (the correct prior incident appears in the top $k$, for
$k\in\{1,5\}$) and mean reciprocal rank (MRR), averaged over queries that have
at least one ground-truth match; on the real dataset we report both the coarse
and strict matching relations of \S\ref{sec:datasets-wol}. \emph{Triage} is a
calibrated classification problem, so beyond accuracy we report the areas under the precision--recall and ROC
curves, and expected calibration error. We also report pages per incident, capturing whether duplicate alerts are
suppressed. Finally,
\emph{cost} is measured per window as wall-clock time and language-model token
usage, relative to a counterfactual cascade that runs every step on every window.

\subsection{Protocol and Implementation}
\label{sec:setup-protocol}

\noindent\emph{Splits.} The synthetic datasets are split by holding out whole
fault families, and the real dataset by holding out two entire project domains
(\S\ref{sec:datasets}); all numbers are reported on the held-out test split.

\noindent\emph{Statistics.} Every reported difference uses a paired bootstrap
($1{,}000$ resamples, $95\%$ confidence intervals) in which the same resampled
windows are scored for both systems, preserving per-window correlation; a
difference is significant when its interval excludes zero. Because we run many comparisons, we control the false-discovery rate across the
test family with the Benjamini--Hochberg procedure and report corrected $q$-values.

\noindent\emph{Where every label comes from.} No label in this paper is produced
by a human annotator, and we state that plainly rather than leaving it to be
inferred. The synthetic triage labels are assigned by construction from the fault
injector: a window inside an injected fault is ticket-worthy, a pre-fault window
is noise, a recovery window is borderline. The real triage labels are derived
mechanically from Jira fields, as described in \S\ref{sec:datasets-wol}. The real
relevance labels are the rule-inferred relations of \S\ref{sec:datasets-wol}. A
language model is used in exactly two places, neither of which produces a
benchmark label: it rewrites synthetic tickets into engineer prose, and it
extracts graph entities offline.

\noindent\emph{Gold validation.} Because the real dataset's matches are inferred
rather than human-labelled, we probe them with an LLM-as-judge. We sample $300$
held-out queries that have a gold match; for each we form two pairs, the query
against a randomly chosen gold match and the query against a memory item sampled
uniformly from outside its gold set, and ask Qwen2.5-7B-Instruct for a single
integer: \emph{``Rate how relevant a past ticket is to a current incident
window\ldots\ 5 = clearly the same or related issue, 1 = unrelated.''} We parse
the first digit of the reply, giving $600$ judgements under greedy decoding at a
fixed seed. We report the gold-versus-control gap, and, because it qualifies the
result, also the fraction of gold matches the judge rates highly
(\S\ref{sec:results-rq4}). The control is sampled uniformly and is therefore an
easy negative; we treat this check as exploratory evidence that the labels carry
signal, not as a substitute for human judgement.

\noindent\emph{Implementation.} The dense retriever is a compact MiniLM
bi-encoder fine-tuned for five epochs with hard negatives mined by lexical search;
the hybrid fuses each retriever's top candidates by reciprocal rank fusion. Graph
entities are extracted from tickets by a small language model and matched by
entity overlap. The retrieval-augmented and judge models run offline with
deterministic (greedy) decoding. Every experiment uses a fixed random seed, and
full configurations and seeds are released with the artifact.

\noindent\emph{Compute environment.} All experiments ran on the Trillium cluster, administered by the Digital
Research Alliance of Canada.\footnote{\url{https://alliancecan.ca}} Each training
and sampling job used a single NVIDIA H100 GPU (80\,GB) with up to 24 CPU cores
and the node's per-GPU share of host memory.
